\documentclass[10pt]{article}
\usepackage[T1]{fontenc}
\usepackage[utf8]{inputenc}
\usepackage[french]{babel}
\usepackage{amsmath,amssymb,mathrsfs,mathpazo}
\usepackage{geometry}
\geometry{a4paper,margin=2cm}
\usepackage{graphicx,array,tabularx,tabularray,colortbl,multirow,multicol,booktabs,enumitem,pifont,ulem,fancybox,diagbox,cancel}
\usepackage{pstricks,pstricks-add,pst-plot,pst-tree,pst-node,pst-text,pst-func,pst-math,pst-bspline,pst-3dplot,pst-eucl,pst-solides3d}
\let\pscalculate\undefined
\usepackage{pst-bezier}
\usepackage{tikz,pgfplots,tkz-tab}
\usetikzlibrary{arrows,arrows.meta,patterns,calc,positioning,trees,decorations.markings,shapes.geometric,shapes.misc,shapes.arrows,intersections,decorations.pathmorphing,decorations.pathreplacing,backgrounds,shadows,babel}
\pgfplotsset{compat=1.18}
\usepgfplotslibrary{fillbetween}
\usepackage[np]{numprint}
\usepackage{listings,makecell,pas-tableur,twoopt}
\let\barre\undefined
\usepackage[pstricks]{tablvar}
\definecolor{darkgreen}{RGB}{0,100,0}
\providecommand{\red}{\color{red}}
\providecommand{\blue}{\color{blue}}
\providecommand{\green}{\color{green}}
\providecommand{\black}{\color{black}}
\providecommand{\white}{\color{white}}
\providecommand{\gray}{\color{gray}}
\providecommand{\cyan}{\color{cyan}}
\providecommand{\magenta}{\color{magenta}}
\providecommand{\yellow}{\color{yellow}}
\providecommand{\violet}{\color{violet}}
\providecommand{\orange}{\color{orange}}
\providecommand{\euro}{\texteuro}
\providecommand{\EUR}{\texteuro}
\providecommand{\ieme}{\textsuperscript{e}}
\providecommand{\degres}{\ensuremath{^\circ}}
\providecommand{\degree}{\ensuremath{^\circ}}
\providecommand{\vv}[1]{\overrightarrow{#1}}
\providecommand{\Coord}[1]{#1}
\providecommand{\decosix}{}
\providecommand{\starredbullet}{\ensuremath{\bullet}}
\pagestyle{empty}
\setlength{\parindent}{0pt}
\usepackage{adjustbox}\DeclareUnicodeCharacter{2236}{:}\DeclareUnicodeCharacter{2190}{\ensuremath{\leftarrow}}
\providecommand{\R}{\mathbb{R}}
\providecommand{\N}{\mathbb{N}}
\providecommand{\Z}{\mathbb{Z}}
\providecommand{\Q}{\mathbb{Q}}
\providecommand{\D}{\mathbb{D}}
\providecommand{\C}{\mathbb{C}}
\providecommand{\e}{{\mathrm{e}}}
\providecommand{\vect}[1]{\overrightarrow{#1}}
\providecommand{\vectt}[1]{\overrightarrow{\mathrm{#1}}}
\providecommand{\vvt}[1]{\overrightarrow{\mathrm{#1}}}
\providecommand{\np}[2][]{{#2~#1}}
\providecommand{\vv}[1]{{\overrightarrow{#1}}}
\renewcommand{\d}{{\mathrm{d}}}
\providecommand{\degres}{{^\circ}}
\providecommand{\degre}{{^\circ}}
\providecommand{\blue}{\color{blue}}
\providecommand{\red}{\color{red}}
\providecommand{\enskip}{\quad}
\providecommand{\llbracket}{\lbrack\!\lbrack}
\providecommand{\rrbracket}{\rbrack\!\rbrack}
\providecommand{\euro}{{\text{€}}}
\providecommand{\no}{{\text{n°}}}
\providecommand{\interval}[2]{{[#1;#2]}}
\providecommand{\intervFO}[2]{{[#1;#2[}}
\providecommand{\dotfill}{\ldots}

\providecommand{\ds}{\displaystyle}
\providecommand{\textding}[1]{#1}
\providecommand{\Oij}{\(\mathrm{O};\vec{\imath},\vec{\jmath}\)}
\providecommand{\Oijk}{\(\mathrm{O};\vec{\imath},\vec{\jmath},\vec{k}\)}
\providecommand{\Ouv}{\(\mathrm{O};\vec{u},\vec{v}\)}
\providecommand{\intervFF}[2]{\ensuremath{[#1;#2]}}
\renewcommand{\euro}{\ifmmode\text{\texteuro}\else\texteuro\fi}
\DeclareUnicodeCharacter{00B0}{\ifmmode^{\circ}\else\textdegree\fi}
\DeclareUnicodeCharacter{20AC}{\ifmmode\text{\texteuro}\else\texteuro\fi}
% Énoncés : Éducation nationale. Transcription et corrigés : APMEP, Denis Vergès et contributeurs. Découpage et classement : CoopMaths, Sébastien Lozano, Éric Elter et contributeurs. Adaptation de la présentation : Maths Explorer.
% CC BY-SA : https://coopmaths.fr/www/about/
% https://coopmaths.fr/alea/static/bac/2022/tex/bac_2022_09_sujet1_metropole_1.tex
\begin{document}
\noindent\textit{APMEP / CoopMaths. CC BY-SA. Présentation adaptée par Maths Explorer.}\par\medskip
\medskip

\emph{Cet exercice est un questionnaire à choix multiples.\\
Pour chacune des questions suivantes, une seule des quatre réponses proposées est exacte. Une réponse fausse, une réponse multiple ou l'absence de réponse à une question ne rapporte ni n'enlève de point.\\
Pour répondre, indiquer sur la copie le numéro de la question et la lettre de la réponse choisie. Aucune justification n'est demandée.}

\medskip

\begin{enumerate}
\item On considère la fonction $g$ définie sur $\R$ par : 

\[ g(x) = \dfrac{2\text{e}^x}{\text{e}^x +1}.\]

\smallskip

La courbe représentative de la fonction $ g$ admet pour asymptote en $+ \infty$ la droite d'équation :

\begin{center}
\begin{tabularx}{\linewidth}{X X X X}
\textbf{a.~} $x = 2$ ;&\textbf{b.~} $y = 2$ ;&\textbf{c.~} $y = 0$ ;&\textbf{d.~} $x = - 1$\\
\end{tabularx}
\end{center}
\item ~

\begin{minipage}{0.38\linewidth}
On considère une fonction $f$ définie et deux fois dérivable sur $\R$.

On appelle $\mathcal{C}$ sa représentation graphique.

On désigne par $f''$ la dérivée seconde de $f$.

On a représenté sur le graphique ci-contre la courbe de $f''$, notée 
$\mathcal{C}''$.
\end{minipage}\hfill
\begin{minipage}{0.58\linewidth}
\begin{center}
\psset{unit=0.8cm,arrowsize=2pt 3}
\begin{pspicture*}(-7,-6.02)(3.02,4)
\psgrid[gridlabels=0pt,subgriddiv=1,gridwidth=0.15pt]
\psaxes[linewidth=1.25pt,Dx=10,Dy=10]{->}(0,0)(-7,-6)(3,4)
\multido{\n=-1+1}{4}{\uput[d](\n,0){\small \n}}
\uput[l](0,1){1}
\psplot[plotpoints=2000,linewidth=1.25pt,linecolor=red]{-7}{2.25}{x 1 add 2.71828 x exp mul x 2 sub mul}
\uput[ul](-5,0.25){\red $\mathcal{C}''$}
\end{pspicture*}
\end{center}
\end{minipage}

\begin{center}
\begin{tabularx}{\linewidth}{*{2}{X}}
\textbf{a.~} $\mathcal{C}$ admet un unique point d'inflexion ; &\textbf{b.~} $f$ est convexe sur l'intervalle $[- 1~;~2 ]$ ; \\
\textbf{c.~} $f$ est convexe sur $]- \infty~;~- 1$] et sur $[2~;~+ \infty[$ ;&\textbf{d.~} $f$ est convexe sur $\R$.\\
\end{tabularx}
\end{center}

\item On donne la suite $\left(u_n\right)$ définie par : $u_0 = 0$ et pour tout entier naturel $n,\: u_{n+1} = \dfrac12 u_n + 1$.

La suite $\left(v_n\right)$, définie pour tout entier naturel $n$ par $v_n = u_n  - 2$, est :

\begin{center}
\begin{tabularx}{\linewidth}{*{2}{X}}
\textbf{a.~} arithmétique de raison $- 2$ ;&\textbf{b.~} géométrique de raison $- 2$ ;\\
\textbf{c.~} arithmétique de raison 1 ;&\textbf{d.~}géométrique de raison  $\dfrac12$.\\
\end{tabularx}
\end{center}

\item  On considère une suite $\left(u_n\right)$ telle que, pour tout entier naturel, on a :

\[1  + \left(\dfrac14\right)^n \leqslant u_n \leqslant 2 - \dfrac{n}{n+1}\]

On peut affirmer que la suite $\left(u_n\right)$ :

\begin{center}
\begin{tabularx}{\linewidth}{*{2}{X}}
\textbf{a.~} converge vers 2 ; &\textbf{b.~} converge vers 1 ;\\
\textbf{c.~} diverge vers $+ \infty$ ;&\textbf{d.~} n'a pas de limite.\\
\end{tabularx}
\end{center}

\item Soit $f$ la fonction définie sur $]0~;~+\infty[$ par $f(x) = x^2 \ln x$. 

Une primitive $F$ de $f$ sur $]0~;~+\infty[$ est définie par :

\begin{center}
\begin{tabularx}{\linewidth}{*{2}{X}}
\textbf{a.~} $F(x) = \dfrac13 x^3 \left(\ln x - \dfrac13 \right)$ ;&\textbf{b.~} $F(x) = \dfrac13 x^3 (\ln x - 1)$ ;\\[8pt]
\textbf{c.~} $F(x) = \dfrac13 x^2$ ;&\textbf{d.~} $F(x) = \dfrac13 x^2 (\ln x - 1)$.\\
\end{tabularx}
\end{center}

\item Pour tout réel $x$, l'expression $2 + \dfrac{3\text{e}^{-x} - 5}{\text{e}^{-x} + 1}$ est égale à : 

\begin{center}
\begin{tabularx}{\linewidth}{*{2}{X}}
\textbf{a.~} $\dfrac{5 - 3\text{e}^x}{1 + \text{e}^x}$ ;&\textbf{b.~} $\dfrac{5 + 3\text{e}^x}{1 - \text{e}^x}$ ;\\[8pt]
\textbf{c.~} $\dfrac{5 + 3\text{e}^x}{1 + \text{e}^x}$ ;&\textbf{d.~} $\dfrac{5 - 3\text{e}^x}{1 - \text{e}^x}$.
\end{tabularx}
\end{center}
\end{enumerate}

\bigskip

\end{document}
